\begin{sheetframe}[50]{cyan}{03｜多変数・変数変換}{領域の形から置換を決める}

\begin{center}
\begin{tikzpicture}[x=1mm,y=1mm]
  \node[labelnode] (shape) at (8,0) {領域の式\\を見る};
  \node[labelnode] (vars) at (39,0) {同じ組合せを\\新しい変数にする};
  \node[labelnode] (inverse) at (75,0) {逆に解いて\\$x=x(u,v),\,y=y(u,v)$};
  \node[labelnode] (range) at (112,0) {不等式を直して\\積分範囲を決める};
  \node[labelnode] (jac) at (148,0) {$|J|$ を掛けて\\積分する};
  \draw[flow] (shape)--(vars);
  \draw[flow] (vars)--(inverse);
  \draw[flow] (inverse)--(range);
  \draw[flow] (range)--(jac);
\end{tikzpicture}
\end{center}

\vspace{.15mm}{\color{line}\rule{\textwidth}{.35pt}}\vspace{.25mm}

\begin{columns}[T,onlytextwidth]
\begin{column}{.305\textwidth}
\subhead{極限・連続・微分可能性}
\begin{center}
\begin{tikzpicture}[x=.82mm,y=.82mm]
  \draw[muted,line width=.4pt] (-18,0)--(20,0) (0,-10)--(0,13);
  \draw[accent,line width=.55pt] (0,0) circle (9);
  \draw[flow] (0,0)--(7,5.4) node[midway,above left,tinylabel] {$r$};
  \draw[accent,line width=.4pt] (4.5,0) arc[start angle=0,end angle=38,radius=4.5];
  \node[tinylabel] at (6.5,1.8) {$\theta$};
  \node[labelnode,anchor=west] at (23,3)
    {$x=r\cos\theta,\ y=r\sin\theta$};
  \node[tinylabel,anchor=west] at (23,-3)
    {$|f(x,y)-L|\le Cr^\alpha,\ \alpha>0\Rightarrow f\to L$};
\end{tikzpicture}
\end{center}
極限が存在しないことは、$y=mx$ や $y=x^2$ など二つの経路で値が違うことを示す。
存在を示すときは経路ではなく、$r$ だけで上から押さえる。
\[
\partial_xf(a),\partial_yf(a)\text{ が存在}
\quad\not\Rightarrow\quad
f\text{ が }a\text{ で微分可能}.
\]
\[
\boxed{
f(a+h)=f(a)+\nabla f(a)\cdot h+o(\norm h)}
\]
$f_x,f_y$ が $a$ の近くで連続なら、上の微分可能性が従う。
\[
D(g\circ f)(a)=Dg(f(a))\,Df(a).
\]

\separator
\subhead{2次 Taylor 展開・停留点}
$f\in C^2$ なら
\[
f(a+h)=f(a)+\nabla f(a)\cdot h
+\frac12h^{\mathsf T}H_f(a)h+o(\norm h^2).
\]
$\nabla f(a)=0$ とし
\[
\Delta=f_{xx}(a)f_{yy}(a)-f_{xy}(a)^2.
\]
\[
\begin{array}{c|c}
\text{条件}&a\text{ の型}\\ \hline
\Delta>0,\ f_{xx}>0&\text{極小}\\
\Delta>0,\ f_{xx}<0&\text{極大}\\
\Delta<0&\text{鞍点}\\
\Delta=0&\text{この判定だけでは不明}
\end{array}
\]

\separator
\subhead{制約つき極値：Lagrange 法}
\[
\boxed{\nabla f=\lambda\nabla g,\qquad g=c.}
\]
例：$f(x,y)=xy$ を $x^2+y^2=1$ 上で調べる。
\[
(y,x)=\lambda(2x,2y)
\Rightarrow y=\pm x.
\]
\[
\max f=\frac12\quad(y=x),\qquad
\min f=-\frac12\quad(y=-x).
\]
\end{column}

\begin{column}{.365\textwidth}
\subhead{変数変換の公式}
$\Phi(u,v)=(x(u,v),y(u,v))$ は $C^1$ で、領域の内部で1対1かつ $J_\Phi\ne0$ とする。
\[
J_\Phi=
\frac{\partial(x,y)}{\partial(u,v)}
=
\det\begin{pmatrix}
x_u&x_v\\y_u&y_v
\end{pmatrix}.
\]
\[
\boxed{
\iint_D f(x,y)\dd x\dd y
=
\iint_{D'}f(\Phi(u,v))\,|J_\Phi|\dd u\dd v.}
\]
ここで $D'$ は $\Phi(D')=D$ となる新しい領域。領域と被積分関数の両方を書き換える。
1対1にならないときは領域を分け、行列式には絶対値を付ける。

\separator
\subhead{楕円：$x=ar\cos\theta,\ y=br\sin\theta$}
$a,b>0$ とする。
\begin{center}
\begin{tikzpicture}[x=.88mm,y=.88mm]
  \draw[accent,line width=.6pt] (-20,0) circle (8);
  \draw[muted,line width=.3pt] (-28,0)--(-12,0) (-20,-8)--(-20,8);
  \node[tinylabel] at (-20,-11) {$0\le r\le1$};
  \draw[flow] (-9,0)--(5,0)
    node[midway,above,tinylabel] {$x=ar\cos\theta$}
    node[midway,below,tinylabel] {$y=br\sin\theta$};
  \draw[accent,line width=.6pt] (24,0) ellipse (15 and 8);
  \draw[muted,line width=.3pt] (9,0)--(39,0) (24,-8)--(24,8);
  \node[tinylabel] at (24,-11)
    {$x^2/a^2+y^2/b^2\le1$};
\end{tikzpicture}
\end{center}
\[
\begin{gathered}
D=\left\{(x,y):\frac{x^2}{a^2}+\frac{y^2}{b^2}\le1\right\},\\
0\le r\le1,\quad0\le\theta<2\pi,\\
|J|=\left|
\det\begin{pmatrix}
a\cos\theta&-ar\sin\theta\\
b\sin\theta&br\cos\theta
\end{pmatrix}\right|=ab\,r.
\end{gathered}
\]
したがって
\[
\boxed{
\iint_D
G\!\left(\frac{x^2}{a^2}+\frac{y^2}{b^2}\right)\dd x\dd y
=2\pi ab\int_0^1G(r^2)\,r\dd r.}
\]
特に
\[
\iint_D1\dd x\dd y=\pi ab,\qquad
\iint_D\left(\frac{x^2}{a^2}+\frac{y^2}{b^2}\right)\dd x\dd y
=\frac{\pi ab}{2}.
\]

\separator
\subhead{円・円環・扇形なら極座標}
\[
x=r\cos\theta,\qquad y=r\sin\theta,\qquad
\dd x\dd y=r\dd r\dd\theta.
\]
\[
R_1^2\le x^2+y^2\le R_2^2
\iff R_1\le r\le R_2.
\]
角度の範囲は図から決める。たとえば第1象限なら
$0\le\theta\le\pi/2$。
\end{column}

\begin{column}{.31\textwidth}
\subhead{$x+y,\ x-y$ が並ぶなら線形変換}
\[
u=x+y,\qquad v=x-y
\]
\[
x=\frac{u+v}{2},\qquad
y=\frac{u-v}{2},\qquad |J|=\frac12.
\]
例：
\[
D=\{(x,y):|x+y|\le1,\ |x-y|\le1\}
\]
は $(u,v)$ 平面の正方形 $[-1,1]^2$ になるので
\[
\operatorname{Area}(D)
=\frac12\int_{-1}^1\!\int_{-1}^1\dd u\dd v=2.
\]

\separator
\subhead{三角形：$s=x+y$ を先に取る}
\begin{center}
\begin{tikzpicture}[x=.8mm,y=.8mm]
  \draw[->,muted,line width=.4pt] (0,0)--(28,0) node[right,tinylabel] {$x$};
  \draw[->,muted,line width=.4pt] (0,0)--(0,22) node[above,tinylabel] {$y$};
  \fill[accent!8] (0,0)--(22,0)--(0,18)--cycle;
  \draw[accent,line width=.6pt] (22,0)--(0,18);
  \node[tinylabel] at (14,12) {$x+y=c$};
  \draw[flow] (0,0)--(12,7) node[midway,above,tinylabel] {$s$};
\end{tikzpicture}
\end{center}
\[
\begin{gathered}
D=\{x\ge0,\ y\ge0,\ x+y\le c\},\\
s=x+y,\qquad t=\frac{x}{x+y},\\
x=st,\qquad y=s(1-t),\\
0\le s\le c,\quad0\le t\le1,\quad |J|=s.
\end{gathered}
\]
$s$ が原点からの大きさ、$t$ が辺上での割合になる。

\separator
\subhead{3次元の球なら球面座標}
\[
\begin{aligned}
x&=r\sin\varphi\cos\theta,\\
y&=r\sin\varphi\sin\theta,\\
z&=r\cos\varphi,
\end{aligned}
\qquad
\dd V=r^2\sin\varphi\dd r\dd\varphi\dd\theta.
\]
球全体では
\[
0\le r\le R,\quad0\le\varphi\le\pi,\quad
0\le\theta<2\pi.
\]

\separator
\subhead{積分順序を逆にするとき}
\[
\int_a^b\int_{\alpha(x)}^{\beta(x)}f(x,y)\dd y\dd x
\]
は、まず領域を図示し、$y$ を固定した横線で左右の端点を読み直す。交点で領域が分かれるなら積分も分ける。
\end{column}
\end{columns}
\end{sheetframe}
